\section{A polynomial bound}\label{sec:cap}
\begin{proposition}\label{lem:cap}
There are real numbers $C_0>0$ and $k\ge0$ such that $H\le C_0R^k$
for every triple.
\end{proposition}
\begin{fact}[{\cite[Fact F1, (2.1)]{Zhang2026Third}}]\label{fact:sy}
For $c>2$ one has $H\le C_1 R^{1/3}(\log R)^3$ with absolute $C_1$.
\end{fact}
\begin{proof}[Proof of Proposition~\ref{lem:cap}]
Apply Fact~\ref{fact:sy} with the radical of the three pairwise coprime sides.
Since $\log R\le R$, $k=20$ and $C_0=\max(1,C_1)$ suffice, including $(1,1,2)$.
\end{proof}

Fix a rational $\eta$ in $(0,1/9)$ and use Table~\ref{tab:parameters}.
Choose the first of $a,b,c$ having least radical and denote it by $s$.
Define $z=c/b,c/a,-a/b$, respectively, when $s=a,b,c$.
These are the ratios in \cite[(3.4)]{Zhang2026Third}, inverted in the first two cases.
For $q\mid s$ they are $q$-units, so inversion preserves $v_q(1-z)=v_q(s)$.
Use the normalizations of \cite[\S2]{Zhang2026Third}, with $v_q(q)=1$.
For $n/d$ in lowest terms with $d>0$ put
$h(n/d)=\log\max(|n|,d)$ and write $h=h(z)$.
Set
\[
  S(s,Y)=\sum_{\substack{q\mid s,\ q\le Y\\q\ {\rm prime}}}v_q(s)\log q .
\]
For the prime $p$ of Section~\ref{sec:sel} put $\ell=\log p$ and $m=v_p(s)$.
The length $t$ in Table~\ref{tab:parameters} is defined in
\eqref{eq:plane-height}. All logarithms are natural.
\parameters
